% GENERATED FILE. Edit canonical lessons, then run npm run generate:books.
% canonical-source-hash: fnv1a64:434e312b4c0d0873
% canonical-lessons: TE-C386-vibhedincu, TE-C386-abhyantaram, TE-C386-aksepincu, TE-C386-vyatirekata, TE-C386-tiraskarincu

\chapter{Disagreeing}
\label{ch:te-disagreeing}
\hlchaptermodality{386}

\begin{chapteropening}
\textbf{By the end of this chapter:} \emph{I can disagree, object and turn things down.}

Complete the last lesson of chapter 386: I can disagree, object and turn things down.
\end{chapteropening}

\section[\texorpdfstring{\te{విభేదించు} (vibhēdiñcu)}{విభేదించు (vibhēdiñcu)}]{\te{విభేదించు} (vibhēdiñcu) — to disagree}

\label{lesson:TE-C386-vibhedincu}



\begin{quote}
Before the new one: say the Telugu for to consent, then the Telugu for backing.
\end{quote}

\subsection*{You'll want to know: \te{విభేదించు}}
\textbf{\te{విభేదించు}} — \emph{vibhēdiñcu} — \textquotedblleft{}to disagree\textquotedblright{}.

\begin{quote}
\te{ఆ} \te{విషయంలో} \te{నేను} \te{మీతో} \te{విభేదిస్తాను}.
\end{quote}

\begin{quote}
\emph{ā viṣayaṁlō nēnu mītō vibhēdistānu.}
\end{quote}

\begin{quote}
On that matter I disagree with you.
\end{quote}

Past (I …): \te{విభేదించాను}, \emph{vibhēdiñcānu}.

Like \te{ఏకీభవించు}, it takes -\te{తో}: \te{మీతో} \te{విభేదిస్తాను}.

\begin{cousinweb}
\textbf{\te{విభేదించు}} (\emph{vibhēdiñcu}) is built on Sanskrit \emph{vibheda}, \textquotedblleft{}a breaking asunder, a division\textquotedblright{} (Monier-Williams), with the Telugu verb ending -\te{ఇంచు}.
\end{cousinweb}

\begin{grammarlens}[title={What you've built}]
One more word for giving reasons and opinions.
\end{grammarlens}

\subsection*{Your turn}
\begin{itemize}
\raggedright
  \item \emph{Say it:} \emph{vibhēdiñcu}
  \item \emph{Say it:} \emph{vibhēdiñcu}, once more
  \item \emph{Say it:} \emph{ā viṣayaṁlō nēnu mītō vibhēdistānu.} — On that matter I disagree with you.
  \item \emph{Recall:} say the Telugu for to consent, then the Telugu for backing, then say \emph{vibhēdiñcu} again
\end{itemize}

\begin{tcolorbox}[breakable,colback=teal!4,colframe=teal!35!black,title={Before you move on}]
What does \te{విభేదించు} mean? (\textquotedblleft{}To disagree\textquotedblright{}.) Say it once more.
\end{tcolorbox}
\section[\texorpdfstring{\te{అభ్యంతరం} (abhyantaraṁ)}{అభ్యంతరం (abhyantaraṁ)}]{\te{అభ్యంతరం} (abhyantaraṁ) — an objection}

\label{lesson:TE-C386-abhyantaram}



\begin{quote}
Before the new one: say the Telugu for backing, then the Telugu for to disagree.
\end{quote}

\subsection*{You'll want to know: \te{అభ్యంతరం}}
\textbf{\te{అభ్యంతరం}} — \emph{abhyantaraṁ} — \textquotedblleft{}an objection\textquotedblright{}.

\begin{quote}
\te{మీకు} \te{ఏమైనా} \te{అభ్యంతరం} \te{ఉందా}?
\end{quote}

\begin{quote}
\emph{mīku ēmainā abhyantaraṁ undā?}
\end{quote}

\begin{quote}
Do you have any objection?
\end{quote}

Plural: \te{అభ్యంతరాలు}, \emph{abhyantarālu}.

A polite opening: \te{మీకు} \te{అభ్యంతరం} \te{లేకపోతే}, if you have no objection.

\begin{cousinweb}
\textbf{\te{అభ్యంతరం}} (\emph{abhyantaraṁ}) is borrowed from Sanskrit \emph{abhyantara}, \textquotedblleft{}interior\textquotedblright{}, and as a noun \textquotedblleft{}the inside\textquotedblright{} (CDIAL 544). In Telugu the word has taken the sense of an objection.
\end{cousinweb}

\begin{grammarlens}[title={What you've built}]
One more word for giving reasons and opinions.
\end{grammarlens}

\subsection*{Your turn}
\begin{itemize}
\raggedright
  \item \emph{Say it:} \emph{abhyantaraṁ}
  \item \emph{Say it:} \emph{abhyantaraṁ}, once more
  \item \emph{Say it:} \emph{mīku ēmainā abhyantaraṁ undā?} — Do you have any objection?
  \item \emph{Recall:} say the Telugu for backing, then the Telugu for to disagree, then say \emph{abhyantaraṁ} again
\end{itemize}

\begin{tcolorbox}[breakable,colback=teal!4,colframe=teal!35!black,title={Before you move on}]
What does \te{అభ్యంతరం} mean? (\textquotedblleft{}An objection\textquotedblright{}.) Say it once more.
\end{tcolorbox}
\section[\texorpdfstring{\te{ఆక్షేపించు} (ākṣēpiñcu)}{ఆక్షేపించు (ākṣēpiñcu)}]{\te{ఆక్షేపించు} (ākṣēpiñcu) — to object; to take exception}

\label{lesson:TE-C386-aksepincu}



\begin{quote}
Before the new one: say the Telugu for to disagree, then the Telugu for an objection.
\end{quote}

\subsection*{You'll want to know: \te{ఆక్షేపించు}}
\textbf{\te{ఆక్షేపించు}} — \emph{ākṣēpiñcu} — \textquotedblleft{}to object; to take exception\textquotedblright{}.

\begin{quote}
\te{కొంత} \te{మంది} \te{ఆ} \te{కొత్త} \te{నియమాన్ని} \te{ఆక్షేపించారు}.
\end{quote}

\begin{quote}
\emph{konta mandi ā kotta niyamānni ākṣēpiñcāru.}
\end{quote}

\begin{quote}
Some people objected to the new rule.
\end{quote}

Past (I …): \te{ఆక్షేపించాను}, \emph{ākṣēpiñcānu}.

Formal and written; in talk people say \te{అభ్యంతరం} \te{చెప్పు}.

\begin{cousinweb}
\textbf{\te{ఆక్షేపించు}} (\emph{ākṣēpiñcu}) is built on Sanskrit \emph{ākṣepa}, \textquotedblleft{}a drawing together; a throwing away\textquotedblright{}, and also \textquotedblleft{}reviling, objection\textquotedblright{} (Monier-Williams), with the Telugu verb ending -\te{ఇంచు}.
\end{cousinweb}

\begin{grammarlens}[title={What you've built}]
One more word for giving reasons and opinions.
\end{grammarlens}

\subsection*{Your turn}
\begin{itemize}
\raggedright
  \item \emph{Say it:} \emph{ākṣēpiñcu}
  \item \emph{Say it:} \emph{ākṣēpiñcu}, once more
  \item \emph{Say it:} \emph{konta mandi ā kotta niyamānni ākṣēpiñcāru.} — Some people objected to the new rule.
  \item \emph{Recall:} say the Telugu for to disagree, then the Telugu for an objection, then say \emph{ākṣēpiñcu} again
\end{itemize}

\begin{tcolorbox}[breakable,colback=teal!4,colframe=teal!35!black,title={Before you move on}]
What does \te{ఆక్షేపించు} mean? (\textquotedblleft{}To object; to take exception\textquotedblright{}.) Say it once more.
\end{tcolorbox}
\section[\texorpdfstring{\te{వ్యతిరేకత} (vyatirēkata)}{వ్యతిరేకత (vyatirēkata)}]{\te{వ్యతిరేకత} (vyatirēkata) — opposition; resistance}

\label{lesson:TE-C386-vyatirekata}



\begin{quote}
Before the new one: say the Telugu for an objection, then the Telugu for to object.
\end{quote}

\subsection*{You'll want to know: \te{వ్యతిరేకత}}
\textbf{\te{వ్యతిరేకత}} — \emph{vyatirēkata} — \textquotedblleft{}opposition; resistance\textquotedblright{}.

\begin{quote}
\te{ఆ} \te{ప్రతిపాదనకు} \te{చాలా} \te{వ్యతిరేకత} \te{వచ్చింది}.
\end{quote}

\begin{quote}
\emph{ā pratipādanaku cālā vyatirēkata vaccindi.}
\end{quote}

\begin{quote}
That proposal met a lot of opposition.
\end{quote}

The noun of \te{వ్యతిరేకించు}, to oppose, which you know; -\te{త} makes an abstract noun.

\begin{cousinweb}
\textbf{\te{వ్యతిరేకత}} (\emph{vyatirēkata}) is built on Sanskrit \emph{vyatireka}, \textquotedblleft{}distinction, separateness; contrast\textquotedblright{} (Monier-Williams), with the abstract ending -\te{త}; the verb \te{వ్యతిరేకించు} comes from the same word.
\end{cousinweb}

\begin{grammarlens}[title={What you've built}]
One more word for giving reasons and opinions.
\end{grammarlens}

\subsection*{Your turn}
\begin{itemize}
\raggedright
  \item \emph{Say it:} \emph{vyatirēkata}
  \item \emph{Say it:} \emph{vyatirēkata}, once more
  \item \emph{Say it:} \emph{ā pratipādanaku cālā vyatirēkata vaccindi.} — That proposal met a lot of opposition.
  \item \emph{Recall:} say the Telugu for an objection, then the Telugu for to object, then say \emph{vyatirēkata} again
\end{itemize}

\begin{tcolorbox}[breakable,colback=teal!4,colframe=teal!35!black,title={Before you move on}]
What does \te{వ్యతిరేకత} mean? (\textquotedblleft{}Opposition; resistance\textquotedblright{}.) Say it once more.
\end{tcolorbox}
\section[\texorpdfstring{\te{తిరస్కరించు} (tiraskariñcu)}{తిరస్కరించు (tiraskariñcu)}]{\te{తిరస్కరించు} (tiraskariñcu) — to reject; to turn down}

\label{lesson:TE-C386-tiraskarincu}



\begin{quote}
Before the new one: say the Telugu for to object, then the Telugu for opposition.
\end{quote}

\subsection*{You'll want to know: \te{తిరస్కరించు}}
\textbf{\te{తిరస్కరించు}} — \emph{tiraskariñcu} — \textquotedblleft{}to reject; to turn down\textquotedblright{}.

\begin{quote}
\te{కంపెనీ} \te{నా} \te{దరఖాస్తును} \te{తిరస్కరించింది}.
\end{quote}

\begin{quote}
\emph{kampenī nā darakhāstunu tiraskariñcindi.}
\end{quote}

\begin{quote}
The company rejected my application.
\end{quote}

Past (I …): \te{తిరస్కరించాను}, \emph{tiraskariñcānu}.

Firm and formal: it goes with applications, offers and ideas.

\begin{cousinweb}
\textbf{\te{తిరస్కరించు}} (\emph{tiraskariñcu}) is built on Sanskrit \emph{tiraskāra}, \textquotedblleft{}a placing aside\textquotedblright{}, and also \textquotedblleft{}disdain\textquotedblright{} (Monier-Williams), with the Telugu verb ending -\te{ఇంచు}.
\end{cousinweb}

\begin{grammarlens}[title={What you've built}]
One more word for giving reasons and opinions.
\end{grammarlens}

\subsection*{Your turn}
\begin{itemize}
\raggedright
  \item \emph{Say it:} \emph{tiraskariñcu}
  \item \emph{Say it:} \emph{tiraskariñcu}, once more
  \item \emph{Say it:} \emph{kampenī nā darakhāstunu tiraskariñcindi.} — The company rejected my application.
  \item \emph{Recall:} say the Telugu for to object, then the Telugu for opposition, then say \emph{tiraskariñcu} again
\end{itemize}

\begin{tcolorbox}[breakable,colback=teal!4,colframe=teal!35!black,title={Before you move on}]
What does \te{తిరస్కరించు} mean? (\textquotedblleft{}To reject; to turn down\textquotedblright{}.) Say it once more.
\end{tcolorbox}
